\section{Fixed-length matrix spanning}\label{sec:lemma2b}

The preceding sections establish the cumulative-spanning part of the
Lemma~2(a) argument from~\cite{Sanz2010Quantum}:
eigenvector spreading gives a cumulative vector span
$K_{D-1}(A,\varphi)=\C^D$.
The next stage is to pass to a span at one fixed length,
namely to find $n_0$ such that $S_{n_0}(A)=\MN{D}$
(i.e.\ word products of \emph{exactly} one length span all matrices).

\begin{definition}[Fixed-length vector span]\label{def:vector_spread_span}
    \lean{MPSTensor.vectorSpreadSpan}
    \leanok
    \uses{def:mps_tensor, def:eval_word}
    The \emph{fixed-length vector span} at length~$n$ is
    \begin{align}
        H_n(A,\varphi)
        &=\spn_{\C}\{A^w\varphi:|w|=n\}\subseteq\C^D.
        \notag
    \end{align}
    This is $S_n(A)\ket{\varphi}$ in the notation
    of~\cite{Sanz2010Quantum}.
\end{definition}

\begin{lemma}[Word span products]\label{lem:word_span_mul}
    \lean{MPSTensor.wordSpan_mul_le}
    \leanok
    \uses{def:word_span}
    One has $S_m(A)S_n(A)\subseteq S_{m+n}(A)$:
    the product of word spans at lengths $m$ and $n$
    is contained in the word span at length $m+n$.
\end{lemma}

\begin{proof}\leanok
    Every product $A^{w_1}A^{w_2}$ with $|w_1|=m$
    and $|w_2|=n$ equals $A^{w_1w_2}$, where
    $|w_1w_2|=m+n$.
\end{proof}

\begin{lemma}[Eigenvector padding]\label{lem:eigenvector_padding}
    \lean{MPSTensor.cumulativeVectorSpan_eq_vectorSpreadSpan_of_eigenvector}
    \leanok
    \uses{def:cumulative_vector_span, def:vector_spread_span}
    Suppose $A^{i_0}\varphi=\mu\varphi$ with $\mu\neq0$.
    Then $K_n(A,\varphi)=H_n(A,\varphi)$ for every~$n$.
    In particular, if $K_n(A,\varphi)=\C^D$, then
    $H_n(A,\varphi)=\C^D$.
\end{lemma}

\begin{proof}
    \leanok
    Any word $w$ with $|w|\le n$ can be padded to exact
    length~$n$ by appending $k=n-|w|$ copies of the index~$i_0$.
    The eigenvector relation gives
    $A^{w\cdot i_0^k}\varphi=\mu^kA^w\varphi$, so the padded vector is a
    nonzero scalar multiple of the original.
    Hence every generator of $K_n$ lies in $H_n$.
    The reverse inclusion $H_n\le K_n$ is immediate.
\end{proof}

\begin{theorem}[Fixed-length vector spanning implies matrix spanning]%
    \label{thm:lemma2b_assembly}
    \lean{MPSTensor.wordSpan_eq_top_of_vectorSpreadSpan_eq_top_of_rankOneBasis}
    \leanok
    \uses{def:word_span, def:vector_spread_span}
    Assume:
    \begin{enumerate}
        \item $H_n(A,\varphi)=\C^D$
              (length-$n$ word products applied to $\varphi$
              span all of~$\C^D$), and
        \item for each standard basis vector $e_j$, the rank-one operator
              $\ketbra{\varphi}{e_j}$ lies in $S_m(A)$.
    \end{enumerate}
    Then $S_{n+m}(A)=\MN{D}$.
\end{theorem}

\begin{proof}
    \leanok
    \uses{lem:word_span_mul}
    Fix a matrix unit $E_{ij}$.
    By hypothesis~(1), there exists $M_i\in S_n(A)$
    with $M_i\varphi=e_i$.
    By hypothesis~(2), $\ketbra{\varphi}{e_j}\in S_m(A)$.
    Their product satisfies
    \begin{align}
        M_i\ketbra{\varphi}{e_j}
        &=\ketbra{M_i\varphi}{e_j}
          =\ketbra{e_i}{e_j}
          =E_{ij}.
        \notag
    \end{align}
    Lemma~\ref{lem:word_span_mul} therefore gives
    $E_{ij}\in S_{n+m}(A)$.
    Since the matrix units span $\MN{D}$,
    $S_{n+m}(A)=\MN{D}$.
\end{proof}

\begin{lemma}[Blocking transfer]\label{lem:blocking_transfer}
    \lean{MPSTensor.wordSpan_eq_top_of_blockTensor_wordSpan_eq_top}
    \lean{MPSTensor.isNBlkInjective_iff_blockTensor_isInjective}
    \leanok
    \uses{def:word_span, def:blocked_tensor}
    If the blocked tensor $A^{[L]}$ satisfies
    $S_n(A^{[L]})=\MN{D}$, then $S_{nL}(A)=\MN{D}$.
    In particular, $A$ is $L$-block injective if and only if $A^{[L]}$ is
    one-site injective.
\end{lemma}

\begin{proof}
    \leanok
    By the blocked-tensor definition (\ref{eq:mps_blocked_tensor}), each
    length-$n$ word in the blocked alphabet
    $\{0,\ldots,d^L-1\}$ decodes to a length-$nL$ word
    in the original alphabet $\{0,\ldots,d-1\}$, so
    $S_n(A^{[L]})\subseteq S_{nL}(A)$.
    For $n=1$, the $d^L$ blocked letters are indexed by the length-$L$ words
    of $A$, giving $S_1(A^{[L]})=S_L(A)$, which proves the stated
    equivalence.
\end{proof}

\begin{theorem}[Bounded injective blocking for left-canonical normal tensors]%
    \label{thm:bounded_injective_blocking_normal_left_canonical}
    \lean{MPSTensor.exists_pos_blockTensor_isInjective_le_pow_four_of_isNormal_leftCanonical}
    \leanok
    \uses{def:normal, def:blocked_tensor, thm:wielandt_bound}
    Let $A$ be a left-canonical normal MPS tensor with bond dimension $D>0$.
    Then there is a positive blocking length $L\le D^4$ such that the blocked
    tensor $A^{[L]}$ is one-site injective.

    This is the blocked-injectivity form of the quantum Wielandt input used in
    the Fundamental Theorem for translation-invariant tensors. It corresponds
    to the statement in
    \cite[Section~II, line~332]{Cirac2016MPDO_arXiv} that every normal tensor
    becomes injective after at most $D^4$ blockings, using the index estimate
    of~\cite[Theorem~1]{Sanz2010Quantum}. The
    left-canonical/trace-preserving hypothesis is included because this is the
    canonical-form context in which the blocked-injectivity input is used.
\end{theorem}

\begin{proof}
    \leanok
    The proof uses the full-Kraus-rank index bound from the quantum Wielandt
    inequality. Since $A$ is normal, the index is attained; the one-step Kraus
    rank is nonzero in the left-canonical case, so the bound
    $(D^2-\dim S_1(A)+1)D^2$ is at most $D^4$. The passage from
    exact word-span fullness to injectivity of the blocked tensor is
    Lemma~\ref{lem:blocking_transfer}.
\end{proof}

\begin{theorem}[Exact quantum-Wielandt blocking length]
    \label{thm:exact_pow_four_block_injectivity_normal_left_canonical}
    \lean{MPSTensor.isNBlkInjective_pow_four_of_isNormal_leftCanonical}
    \leanok
    \uses{thm:bounded_injective_blocking_normal_left_canonical,
        lem:blocking_transfer,
        thm:wordspan_blk_inj_succ}
    Let $A$ be a left-canonical normal MPS tensor with bond dimension $D>0$.
    Then $S_{D^4}(A)=\MN{D}$; equivalently, $A$ is
    $D^4$-block injective.
\end{theorem}

\begin{proof}
    \leanok
    \uses{thm:bounded_injective_blocking_normal_left_canonical,
        lem:blocking_transfer,
        thm:wordspan_blk_inj_succ}
    Choose $L$ with $0<L\le D^4$ and $S_L(A)=\MN{D}$, and write
    $D^4=L+k$.
    Induction on $k$, using positive-length propagation at each step, gives
    $S_{L+k}(A)=\MN{D}$.
\end{proof}

\begin{theorem}[Common exact block-injectivity length]
    \label{thm:common_exact_block_injectivity_normal_left_canonical}
    \lean{MPSTensor.%
        exists_common_isNBlkInjective_of_isNormal_leftCanonical}
    \leanok
    \uses{def:l_blk_injective, def:normal,
        thm:bounded_injective_blocking_normal_left_canonical}
    Let $(A_j)_{j=1}^g$ be a finite family of positive-bond-dimension,
    left-canonical normal tensors. There is a positive integer $L$ such that
    every $A_j$ is $L$-block injective.
\end{theorem}

\begin{proof}\leanok
    Apply
    Theorem~\ref{thm:bounded_injective_blocking_normal_left_canonical} to each
    block. The product of the finitely many positive lengths is a common
    positive multiple, and exact block injectivity persists under positive
    multiples.
\end{proof}

\begin{corollary}[Uniform injective blocking for separated normal-canonical families]%
    \label{cor:uniform_injective_blocking_normal_bnt}
    \lean{MPSTensor.IsNormalCanonicalFormBNT.exists_common_blockTensor_isInjective,
        MPSTensor.exists_common_blockTensor_isInjective_two_of_isNormalCanonicalFormBNT}
    \leanok
    \uses{thm:normal_tp_primitive_irred,
        thm:bounded_injective_blocking_normal_left_canonical}
    Let $(A_x)_{x\in X}$ be a finite separated normal-canonical family whose
    blocks have positive virtual bond dimension. Then there is a positive
    integer $L$ such that every blocked tensor $A_x^{[L]}$ is one-site
    injective. The same conclusion holds simultaneously for two finite
    separated normal-canonical families whose blocks have positive virtual bond
    dimension.

    Here the family is a basis of normal-canonical representatives in the sense
    of Definition~\ref{def:normal_cf_bnt}; the blocks are distinct and the weight
    moduli are non-increasing but need not be strictly ordered. It assumes that
    the representative family has already been selected; it does not reconstruct
    the full CPSV sector decomposition with repeated copies inside one sector.
    The restriction is recorded in
    \path{docs/paper-gaps/ft_one_copy_scope_restriction.tex}.
\end{corollary}

\begin{proof}
    \leanok
    Apply Theorem~\ref{thm:bounded_injective_blocking_normal_left_canonical}
    to each block. Taking the product of the finitely many resulting positive
    blocking lengths gives a common positive multiple, and injectivity persists
    under positive further blocking.
\end{proof}

\begin{remark}[Relation with the rank-one extraction lemma]\label{rem:lemma2b_status}\leanok
    In~\cite{Sanz2010Quantum}, Lemma~2(b) combines a rank-one
    extraction for fixed-length spans with a product
    argument turning fixed-length vector spanning into fixed-length
    matrix spanning.
    Theorem~\ref{thm:lemma2b_assembly} is the second step.
    Theorem~\ref{thm:rankone_extraction_blockTensor} supplies the
    rank-one hypothesis in blocked form, and
    Theorem~\ref{thm:wielandt_lemma2b} then gives a blocked,
    fixed-length word-span saturation result.
\end{remark}

\begin{remark}[Use of the quantum Wielandt theorem in the MPS route]
    The source papers use the quantum Wielandt theorem as an input converting
    normality into injectivity after blocking. In the notation of this chapter,
    the input has two distinct conclusions. First,
    Theorem~\ref{thm:wielandt_bound} gives cumulative saturation
    $T_{D^2}(A)=\MN{D}$. This alone is not the injectivity statement used in
    canonical form, because injectivity requires products of one fixed length.
    Second, Lemma~2(b) of~\cite{Sanz2010Quantum}, in the form of
    Theorem~\ref{thm:wielandt_lemma2b}, supplies a length $N$ for which
    $S_N(A)=\MN{D}$. This is the statement that matches the injectivity of
    $\Gamma_N$ and can be transported through blocking by
    Lemma~\ref{lem:blocking_transfer}. Thus any later use of Wielandt in the
    Fundamental Theorem must specify whether it is using cumulative span,
    exact-length span, or the blocked version of the exact-length conclusion.
\end{remark}

\section{Blocked tensor and rectangular span}\label{sec:rectangular_span}

This section extends the Lemma~2(b) argument to the
\emph{blocked tensor}~$A^{[L]}$ and develops the ``rectangular span''
used to produce rank-one elements.

\begin{theorem}[Blocking preserves normality]\label{thm:isNormal_blockTensor}
    \lean{MPSTensor.isNormal_blockTensor_of_isNormal}
    \leanok
    \uses{def:normal, def:blocked_tensor}
    If $A$ is normal and $L>0$, then the blocked tensor
    $A^{[L]}$ is also normal.
\end{theorem}

\begin{proof}\leanok
    \uses{lem:word_span_le_blockTensor, lem:word_span_mul}
    Choose $N$ with $S_N(A)=\MN{D}$. Then $\Id\in S_N(A)$.
    For any $M\in S_N(A)$, write
    $\Id=\sum_r c_rA^{w_r}$ with each $|w_r|=N$. Then
    \begin{align}
        M
        &=M\Id
          =\sum_r c_rMA^{w_r}
          \in S_N(A)S_N(A)
          \subseteq S_{2N}(A).
        \label{eq:wld_blocknormal_self_inclusion}
    \end{align}
    By Lemma~\ref{lem:word_span_mul}, iterating
    (\ref{eq:wld_blocknormal_self_inclusion}) gives
    $S_N(A)\subseteq S_{kN}(A)$ for every $k\ge1$.
    Taking $k=L$ yields
    $\MN{D}=S_N(A)\subseteq S_{NL}(A)$.
    By Lemma~\ref{lem:word_span_le_blockTensor},
    $S_{NL}(A)\subseteq S_N(A^{[L]})$, so
    $S_N(A^{[L]})=\MN{D}$.
\end{proof}

\begin{lemma}[Chunking: word span embeds into blocked word span]%
    \label{lem:word_span_le_blockTensor}
    \lean{MPSTensor.wordSpan_le_wordSpan_blockTensor}
    \leanok
    \uses{def:word_span, def:blocked_tensor}
    One has $S_{nL}(A)\subseteq S_n(A^{[L]})$ for all $n,L$.
\end{lemma}

\begin{proof}\leanok
    \uses{lem:word_span_mul}
    Every length-$nL$ word in the original alphabet
    can be split into $n$ consecutive chunks of length~$L$,
    each of which is a single blocked Kraus operator by
    (\ref{eq:mps_blocked_tensor}).
    The claim follows by induction on~$n$, using
    Lemma~\ref{lem:word_span_mul}.
\end{proof}

\begin{theorem}[Word eigenvector gives blocked eigenvector]%
    \label{thm:blockTensor_single_eigenvector}
    \lean{MPSTensor.blockTensor_single_eigenvector}
    \leanok
    \uses{def:blocked_tensor}
    If $A^{w_0}\varphi=\mu\varphi$ for a word $w_0$
    of length~$L$, then $(A^{[L]})^{j_0}\varphi=\mu\varphi$,
    where $j_0$ is the encoded blocked index corresponding
    to $w_0$.
\end{theorem}

\begin{proof}\leanok
    By (\ref{eq:mps_blocked_tensor}),
    $(A^{[L]})^{j_0}=A^{w_0}$.
\end{proof}

\begin{theorem}[Blocked fixed-length matrix spanning]\label{thm:wielandt_blocked_assembly}
    \lean{MPSTensor.wielandt_blocked_assembly}
    \leanok
    \uses{def:word_span, def:normal, def:blocked_tensor}
    Suppose $A$ is normal and $L>0$, and we have:
    \begin{enumerate}
        \item a word $w_0$ of length $L$ with eigenvector
              $A^{w_0}\varphi=\mu\varphi$
              ($\mu\neq0$, $\varphi\neq0$),
        \item a word $\tau_0$ of length $L$ with transpose
              eigenvector
              $(A^{\tau_0})^T\psi=\nu\psi$
              ($\nu\neq0$, $\psi\neq0$),
        \item a rank-one element
              $\varphi\psi^T\in S_m(A^{[L]})$
              for some $m$.
    \end{enumerate}
    Then there exists $N$ such that $S_N(A)=\MN{D}$.

    \noindent\emph{Remark.}
    The proof gives the explicit bound
    $N=((D{-}1)+m+(D{-}1))L$; the theorem states only the
    existential conclusion.
\end{theorem}

\begin{proof}
    \leanok
    \uses{thm:isNormal_blockTensor, thm:blockTensor_single_eigenvector,
        thm:eigenvector_spreading, lem:blocking_transfer,
        thm:lemma2b_assembly}
    Set $B=A^{[L]}$, which is normal by
    Theorem~\ref{thm:isNormal_blockTensor}. The word eigenvectors transfer to
    single-index eigenvectors of $B$ and $B^T$ by
    Theorem~\ref{thm:blockTensor_single_eigenvector}. Eigenvector spreading
    and padding then give
    $H_{D-1}(B,\varphi)=H_{D-1}(B^T,\psi)=\C^D$.

    For each standard basis vector $e_j$, choose
    $R_j\in S_{D-1}(B^T)$ with $R_j\psi=e_j$. Reversing each word after
    transposition shows that $R_j^T\in S_{D-1}(B)$, and therefore
    $(\varphi\psi^T)R_j^T=\varphi e_j^T$ belongs to
    $S_{m+D-1}(B)$. Theorem~\ref{thm:lemma2b_assembly}, applied with
    $n=D-1$ and $m'=m+D-1$, now gives
    $S_{(D-1)+m+(D-1)}(B)=\MN{D}$. Finally,
    Lemma~\ref{lem:blocking_transfer} transfers this identity to
    $S_{((D-1)+m+(D-1))L}(A)=\MN{D}$.
\end{proof}

\begin{theorem}[Rank-one extraction for blocked tensor]%
    \label{thm:rankone_extraction_blockTensor}
    \lean{MPSTensor.exists_rankOne_mem_wordSpan_blockTensor}
    \leanok
    \uses{def:word_span, def:normal, def:blocked_tensor}
    Suppose $A$ is normal and $N_0\ge1$ with
    $S_{N_0}(A)=\MN{D}$.
    Then there exist words $\sigma_0,\tau_0$ of length~$N_0$,
    nonzero vectors $\varphi,\psi$, nonzero scalars $\mu,\nu$,
    and $m\in\N$ such that:
    \begin{enumerate}
        \item $A^{\sigma_0}\varphi=\mu\varphi$,
        \item $(A^{\tau_0})^T\psi=\nu\psi$,
        \item $\varphi\psi^T\in S_m(A^{[N_0]})$.
    \end{enumerate}
\end{theorem}

\begin{proof}
    \leanok
    \uses{def:blocked_tensor, thm:isNormal_blockTensor,
        thm:eigenvector_from_trace}
    Set $B=A^{[N_0]}$.
    By (\ref{eq:mps_blocked_tensor}) and
    $S_{N_0}(A)=\MN{D}$, one has $S_1(B)=\MN{D}$.
    Hence $\Id\in S_1(B)$, so in the nontrivial case $D\ge1$
    not all blocked generators can have trace zero.
    Choose $j_0$ with $\tr(B^{j_0})\neq0$.
    By Theorem~\ref{thm:eigenvector_from_trace},
    $B^{j_0}$ has a nonzero eigenvector $\varphi$
    with eigenvalue $\mu\neq0$.
    Writing $j_0$ as the blocked index corresponding to a word
    $\sigma_0$ of length~$N_0$ gives
    $A^{\sigma_0}\varphi=\mu\varphi$.
    Applying the same argument to the transpose blocked tensor gives
    a word $\tau_0$ of length~$N_0$, a nonzero vector $\psi$,
    and $\nu\neq0$ with $(A^{\tau_0})^T\psi=\nu\psi$.
    Since $A$ is normal and $N_0\ge1$,
    the blocked tensor $B$ is normal by
    Theorem~\ref{thm:isNormal_blockTensor}.
    Therefore there exists $M$ with $S_M(B)=\MN{D}$.
    In particular, $\varphi\psi^T\in S_M(B)$.
    Taking $m=M$ proves~(3).
\end{proof}

\begin{theorem}[Fixed-length word-span saturation via blocking]\label{thm:wielandt_lemma2b}
    \lean{MPSTensor.wielandt_lemma2b}
    \leanok
    \uses{def:word_span, def:normal}
    If $A$ is a normal MPS tensor with bond dimension~$D$,
    then there exists $N$ such that $S_N(A)=\MN{D}$.
    This gives the fixed-length spanning conclusion corresponding to
    \cite[Lemma~2(b)]{Sanz2010Quantum}.
\end{theorem}

\begin{proof}
    \leanok
    \uses{thm:rankone_extraction_blockTensor,
        thm:wielandt_blocked_assembly}
    By normality, some $N_0$ has $S_{N_0}(A)=\MN{D}$.
    If $N_0=0$, we are done.
    Otherwise, Theorem~\ref{thm:rankone_extraction_blockTensor}
    produces eigenvectors and a rank-one element
    $\varphi\psi^T\in S_m(A^{[N_0]})$.
    Theorem~\ref{thm:wielandt_blocked_assembly} then gives
    $S_{((D-1)+m+(D-1))N_0}(A)=\MN{D}$.
\end{proof}

\begin{remark}[Blocking argument for fixed-length spanning]%
    \label{rem:blocked_lemma2b_gap}\leanok
    Theorem~\ref{thm:wielandt_lemma2b} is obtained by first extracting a
    rank-one operator in a blocked tensor
    (Theorem~\ref{thm:rankone_extraction_blockTensor}) and then
    applying the fixed-length matrix-spanning theorem
    (Theorem~\ref{thm:wielandt_blocked_assembly}).
    Thus the fixed-length span of $A$ is reached through a blocked
    auxiliary tensor rather than by a direct unblocked argument.
\end{remark}

\section{Primitive MPS tensors}

\begin{definition}[Primitive MPS tensor]\label{def:primitive_mps}
    \lean{MPSTensor.IsPrimitiveMPS}
    \leanok
    \uses{def:mps_tensor, def:transfer_map, def:fixed_point_proj}
    Assume $D>0$.
    A tensor $A$ together with a matrix $\rho\in\MN{D}$ is
    \emph{primitive} if
    \begin{align}
        \sum_i(A^i)^\dagger A^i
        &=\Id, \notag\\
        \rho
        &\ge0,\qquad \rho\neq0,\qquad \E_A(\rho)=\rho, \notag\\
        \rho_{\spec}(\E_A-P)
        &<1,
        \notag
    \end{align}
    where $P$ is the fixed-point projection associated to $\rho$.
    When the choice of $\rho$ is irrelevant, we simply say that $A$
    is a primitive MPS tensor.
    This is weaker than the primitive condition of~\cite{Sanz2010Quantum},
    which requires $\rho>0$; see Remark~\ref{rem:prim_definition_mismatch}.
    The stronger positive-definite hypothesis is imposed separately in
    Lemma~\ref{lem:primitive_implies_irreducible}
    and Theorem~\ref{thm:normal_from_primitive_posdef}.
\end{definition}

\begin{theorem}[Primitive overlap convergence]\label{thm:prim_overlap_one}
    \lean{MPSTensor.IsPrimitiveMPS.overlap_tendsto_one}
    \leanok
    \uses{def:primitive_mps, def:mpv_overlap}
    If $A$ is a primitive MPS tensor, then $O_{AA}(N)\to1$.
\end{theorem}

\begin{proof}\leanok\uses{thm:primitive_overlap}
    The claim follows from the self-overlap convergence under a
    complementary transfer-map gap
    (Theorem~\ref{thm:primitive_overlap}).
\end{proof}

\begin{theorem}[Peripheral primitivity and irreducibility give normalized overlap]
    \label{thm:peripheral_primitive_irreducible_overlap_one}
    \lean{MPSTensor.%
        overlap_tendsto_one_of_peripheralPrimitive_of_irreducible}
    \leanok
    \uses{def:primitive_channel, def:irreducible_tensor, def:mpv_overlap}
    Let $A$ have positive bond dimension and satisfy
    $\sum_i(A^i)^\dagger A^i=\Id$.
    If $A$ is tensor-irreducible and the transfer map is primitive, then
    $O_{AA}(N)\to1$.
\end{theorem}

\begin{proof}\leanok\uses{thm:prim_overlap_one}
    Irreducible Perron--Frobenius theory supplies a nonzero positive fixed
    point. Peripheral primitivity gives a strict spectral-radius bound on the
    complementary transfer map. Thus $A$ is a primitive MPS tensor, and
    Theorem~\ref{thm:prim_overlap_one} applies.
\end{proof}

\section{Primitivity and normality}

Throughout this section, primitivity means the condition of
Definition~\ref{def:primitive_mps}: the tensor is in TP gauge,
it has a nonzero PSD fixed point, and the complementary map has
spectral radius $<1$.

\begin{lemma}[Primitive MPS tensors have nonzero word products]\label{lem:prim_nonzero_word}
    \lean{MPSTensor.exists_nonzero_evalWord_of_isPrimitiveMPS}
    \leanok
    \uses{def:primitive_mps, def:eval_word}
    If $A$ is a primitive MPS tensor,
    then for every $n\ge0$ there exists a word $w$ of length
    \emph{exactly}~$n$ with $A^w\neq0$.
\end{lemma}

\begin{proof}\leanok
    The PSD fixed point $\rho$ of $\E_A$ satisfies
    $\E_A^n(\rho)=\rho\neq0$.
    Since
    $\E_A^n(\rho)=\sum_{|w|=n}A^w\rho(A^w)^\dagger$,
    at least one summand is nonzero, hence some $A^w\neq0$.
\end{proof}

\begin{lemma}[Iterated transfer does not vanish]\label{lem:transfer_pow_ne_zero}
    \lean{MPSTensor.transferMap_pow_ne_zero_of_isPrimitiveMPS}
    \leanok
    \uses{def:primitive_mps, def:transfer_map}
    If $A$ is a primitive MPS tensor, then
    $\E_A^n\neq0$ for all $n\ge0$.
\end{lemma}

\begin{proof}\leanok\uses{lem:prim_nonzero_word}
    By Lemma~\ref{lem:prim_nonzero_word}, there exists a word~$w$ of
    length~$n$ with $A^w\neq0$.
    The expansion
    $\E_A^n(\Id)=\sum_{|u|=n}A^u(A^u)^\dagger$
    contains the nonzero positive semidefinite summand
    $A^w(A^w)^\dagger$, so $\E_A^n(\Id)\neq0$.
\end{proof}

\subsection{From primitivity to normality}\label{subsec:prim_to_normal}

The primitivity condition of Definition~\ref{def:primitive_mps}
combines TP normalization, a nonzero PSD fixed point, and a complementary
transfer-map gap. It is connected to normality through the following results.

\begin{lemma}[Fixed-point uniqueness from a complementary transfer-map gap]%
    \label{lem:fixedpoint_unique_spectral}
    \lean{MPSTensor.IsPrimitiveMPS.fixedPoint_unique}
    \leanok
    \uses{def:primitive_mps}
    If $A$ is a primitive MPS tensor
    with PSD fixed point $\rho$,
    then every fixed point $\sigma$ of $\E_A$ is proportional
    to $\rho$: $\sigma=\frac{\tr(\sigma)}{\tr(\rho)}\rho$.
\end{lemma}

\begin{proof}\leanok
    Set $\sigma'=\sigma-\frac{\tr(\sigma)}{\tr(\rho)}\rho$.
    Then $\tr(\sigma')=0$ and $(\E_A-P_\rho)(\sigma')=\sigma'$.
    If $\sigma'\neq0$, then $1$ is an eigenvalue of $\E_A-P_\rho$,
    contradicting the complementary transfer-map gap
    $\rho_{\spec}(\E_A-P_\rho)<1$.
\end{proof}

\begin{lemma}[Complement powers tend to zero]%
    \label{lem:complement_pow_zero}
    \lean{MPSTensor.IsPrimitiveMPS.complement_pow_tendsto_zero}
    \leanok
    \uses{def:primitive_mps}
    If $A$ is a primitive MPS tensor with PSD fixed point $\rho$,
    then $(\E_A-P_\rho)^n\to0$ in operator norm.
\end{lemma}

\begin{proof}\leanok
    The spectral radius of $\E_A-P_\rho$ is strictly less
    than $1$ by the complementary transfer-map gap hypothesis.
    Apply the general result that powers of an operator
    with spectral radius $<1$ tend to zero.
\end{proof}

\begin{remark}[Primitivity: transfer-map gap vs.\ positive definiteness]%
    \label{rem:prim_definition_mismatch}\leanok
    The primitivity condition used in Section~\ref{subsec:prim_to_normal}
    has three parts: TP normalization $\sum_i(A^i)^\dagger A^i=\Id$,
    a nonzero PSD fixed point $\rho$, and the complementary transfer-map gap
    $\rho_{\spec}(\E_A-P_\rho)<1$.
    The notion of ``primitive'' in
    \cite[Proposition~3]{Sanz2010Quantum} additionally requires
    $\rho$ to be positive definite (full rank).
    Without positive definiteness, the transfer-map gap condition alone
    does \emph{not} imply irreducibility or normality.
    A counterexample is $D=2$, $\rho=\diag(1,0)$.

    With the additional hypothesis $\rho>0$,
    Lemma~\ref{lem:primitive_implies_irreducible} gives
    irreducibility.
    Together with the Burnside argument below
    (Section~\ref{subsec:burnside_full_span})
    and Theorem~\ref{thm:algspan_to_normal}, one obtains
    \begin{align}
        \text{primitive}+\rho\text{ positive definite}
        &\Rightarrow\text{irreducible}
          \Rightarrow\algspn(A)=\MN{D}
          \Rightarrow\text{normal}.
        \notag
    \end{align}
    The last step uses the aperiodicity condition
    $\Id\in S_1(A)$ from Remark~\ref{rem:aperiodicity}.
    After restoring the positive-definite hypothesis,
    \cite[Proposition~3]{Sanz2010Quantum} identifies that
    primitive/strongly irreducible condition with eventual full
    Kraus rank, i.e.\ with our notion of normality.
\end{remark}

\subsection{Primitive implies irreducible}\label{subsec:prim_implies_irred}

\begin{lemma}[Primitive tensors with positive-definite fixed point are irreducible]%
    \label{lem:primitive_implies_irreducible}
    \lean{MPSTensor.isIrreducibleTensor_of_isPrimitiveMPS_of_posDef}
    \leanok
    \uses{def:primitive_mps, def:irreducible_tensor}
    If $A$ is a primitive MPS tensor with PSD fixed point $\rho$
    and $\rho$ is positive definite, then $A$ is an irreducible tensor.
    This is part of~\cite[Proposition~3]{Sanz2010Quantum}
    (see also \cite[Section~2.3]{Cirac2016MPDO_arXiv}).
\end{lemma}

\begin{proof}
    \leanok
    \uses{def:has_inv_proj, lem:complement_pow_zero}
    By contrapositive.
    Assume $A$ has an invariant projection~$P$
    ($P\neq0$, $P\neq\Id$, $(\Id-P)A^iP=0$ for all~$i$).
    Set $\sigma=P\rho P$.
    \begin{enumerate}
        \item \emph{Invariance of $\sigma$ under $\E_A^n$:}
              the projection relation
              $(\Id-P)A^iP=0$ implies
              $(\Id-P)\E_A^n(\sigma)=0$ for all~$n$
              (by induction on~$n$).
        \item \emph{Convergence:}
              By Lemma~\ref{lem:complement_pow_zero},
              $\E_A^n(\sigma)\to\frac{\tr(\sigma)}{\tr(\rho)}\rho$.
              Hence
              $(\Id-P)\frac{\tr(\sigma)}{\tr(\rho)}\rho=0$.
        \item \emph{Nonzero scalar:}
              Since $P\neq0$ and $\rho$ is positive definite,
              $\tr(P\rho P)>0$, so
              $\frac{\tr(\sigma)}{\tr(\rho)}\neq0$.
              Therefore $(\Id-P)\rho=0$.
        \item \emph{Contradiction:}
              Since $\rho$ is positive definite, it is a unit
              in $\MN{D}$.
              From $(\Id-P)\rho=0$ we get $\Id-P=0$,
              i.e.\ $P=\Id$, contradicting $P\neq\Id$.
    \end{enumerate}
\end{proof}

\subsection{From irreducibility to full spanning via Burnside's theorem}
\label{subsec:burnside_full_span}

The irreducibility condition (Definition~\ref{def:irreducible_tensor}) is formulated
in terms of invariant \emph{projections}.
For Burnside-style arguments it is more natural to work with invariant
\emph{subspaces} of $\C^D$ under the Kraus operators.

\begin{definition}[Algebra span]\label{def:alg_span}
    \lean{MPSTensor.algSpan}\leanok
    \uses{def:mps_tensor}
    For an MPS tensor $A$, let $\algspn(A)$ be the unital
    $\C$-subalgebra of $\MN{D}$ generated by the matrices
    $\{A^i\}_{i=0}^{d-1}$:
    $\algspn(A)=\C\langle A^i:i=0,\ldots,d-1\rangle\subseteq\MN{D}$.
\end{definition}

\begin{definition}[Invariant submodule]\label{def:invariant_submodule}
    \lean{MPSTensor.IsInvariantSubmodule}\leanok
    \uses{def:mps_tensor}
    A subspace $W\le\C^D$ is \emph{$A$-invariant} if
    $A^iW\subseteq W$ for every~$i$.
\end{definition}

\begin{definition}[Irreducible action]\label{def:irreducible_action}
    \lean{MPSTensor.IsIrreducibleAction}\leanok
    \uses{def:invariant_submodule}
    The Kraus operators of $A$ act \emph{irreducibly} on $\C^D$ if every
    $A$-invariant subspace is trivial, i.e.\ equal to $0$ or to $\C^D$.
\end{definition}

\begin{theorem}[Irreducible tensor implies irreducible action]\label{thm:irreducible_tensor_to_action}
    \lean{MPSTensor.isIrreducibleAction_of_isIrreducibleTensor}\leanok
    \uses{def:irreducible_tensor, def:irreducible_action}
    If $A$ is irreducible as a tensor (Definition~\ref{def:irreducible_tensor}),
    then the Kraus operators act irreducibly on $\C^D$
    (Definition~\ref{def:irreducible_action}).
\end{theorem}

\begin{proof}\leanok
    Let $W\le\C^D$ be an $A$-invariant subspace.
    In finite dimension the orthogonal projection $P$ onto $W$ exists.
    Invariance of $W$ implies $(\Id-P)A^iP=0$ for all~$i$.
    Thus $P$ is a nontrivial invariant projection unless $W=0$ or
    $W=\C^D$, contradicting irreducibility of the tensor.
\end{proof}

\begin{theorem}[Burnside's theorem for Kraus algebras]\label{thm:burnside_matrix}
    \lean{MPSTensor.burnside_matrix}\leanok
    \uses{def:alg_span, def:irreducible_action}
    Assume $D>0$.
    If $A$ acts irreducibly on $\C^D$, then the generated unital subalgebra is
    the full matrix algebra: $\algspn(A)=\MN{D}$.
    This is the complex finite-dimensional case of Burnside's theorem
    (equivalently, Jacobson's density theorem);
    see~\cite{Jacobson2009BasicAlgebraII}.
\end{theorem}

\begin{proof}\leanok
    Transport $\algspn(A)$ along the algebra equivalence
    $\MN{D}\simeq_{\C}\End_{\C}(\C^D)$.
    The irreducible-action hypothesis makes $\C^D$ a simple module for this
    algebra.
    Over the algebraically closed field $\C$, Schur's lemma identifies the
    endomorphisms commuting with the action with scalars, and Jacobson density
    then forces the action map to be surjective onto all of
    $\End_{\C}(\C^D)$.
\end{proof}

\begin{remark}\leanok
    This Burnside/Jacobson-density step is a substantial external algebraic
    ingredient in the chapter. We use it here as the finite-dimensional complex
    Burnside theorem, namely the statement that an irreducible matrix algebra
    action already generates the full endomorphism algebra.
\end{remark}

\subsection{From primitivity to strong irreducibility and normality}%
\label{subsec:prim_normality_connection}

The results in this subsection remove the aperiodicity assumption
from the normality conclusions.
They connect primitivity in Definition~\ref{def:primitive_mps},
formulated via a complementary transfer-map gap, directly to normality
(eventually full Kraus rank),
using the convergence of the transfer-map iterates to the
projection onto the fixed-point subspace.

\begin{lemma}[Primitive tensors with positive-definite fixed point have irreducible transfer map]%
    \label{lem:irreducible_map_from_primitive}
    \lean{MPSTensor.isIrreducibleMap_of_isPrimitiveMPS_of_posDef}
    \leanok
    \uses{def:primitive_mps}
    If $A$ is a primitive MPS tensor with positive-definite fixed
    point~$\rho$, then the transfer map $\E_A$ is irreducible.

    The proof uses fixed-point uniqueness
    (Lemma~\ref{lem:fixedpoint_unique_spectral}):
    every PSD fixed point is proportional to $\rho$,
    so Wolf's criterion for irreducibility
    (positive-definite fixed point implies irreducibility)
    applies directly.
\end{lemma}

\begin{proof}
    \leanok
    \uses{lem:fixedpoint_unique_spectral}
    By Lemma~\ref{lem:fixedpoint_unique_spectral},
    if $\sigma\ge0$ and $\E_A(\sigma)=\sigma$, then
    $\sigma=\frac{\tr(\sigma)}{\tr(\rho)}\rho$.
    Since $\rho>0$, Wolf's uniqueness criterion for
    irreducibility applies.
\end{proof}

\begin{lemma}[Primitive tensors with positive-definite fixed point are strongly irreducible]%
    \label{lem:strongly_irred_from_primitive}
    \lean{MPSTensor.isStronglyIrreduciblePaper_of_isPrimitiveMPS_of_posDef}
    \leanok
    \uses{def:primitive_mps}
    If $A$ is a primitive MPS tensor with positive-definite fixed
    point~$\rho$, then $A$ is strongly irreducible in the sense of~%
    \cite[Proposition~3(c)]{Sanz2010Quantum}:
    $\E_A$ is irreducible, the fixed point $\rho$ is positive
    definite, and the peripheral spectrum of $\E_A$ is $\{1\}$.
\end{lemma}

\begin{proof}
    \leanok
    \uses{lem:irreducible_map_from_primitive, thm:prim_overlap_one}
    Combine Lemma~\ref{lem:irreducible_map_from_primitive}
    (irreducibility) with the peripheral-spectrum primitivity from
    the transfer-map gap hypothesis
    (Theorem~\ref{thm:prim_overlap_one}).
\end{proof}

\begin{theorem}[Primitive MPS with positive-definite fixed point is normal]%
    \label{thm:normal_from_primitive_posdef}
    \lean{MPSTensor.isNormal_of_isPrimitiveMPS_with_posDef}
    \leanok
    \uses{def:primitive_mps, def:normal}
    If $A$ is a primitive MPS tensor with positive-definite fixed
    point~$\rho$, then $A$ is normal
    (has eventually full Kraus rank).
    No aperiodicity assumption is needed.

    This is the key connection between primitivity in
    Definition~\ref{def:primitive_mps}, formulated via a complementary
    transfer-map gap, and normality; it removes the need for the
    aperiodicity-based argument.
\end{theorem}

\begin{proof}
    \leanok
    \uses{lem:strongly_irred_from_primitive}
    By Lemma~\ref{lem:strongly_irred_from_primitive},
    $A$ is strongly irreducible. The standard equivalence in
    \cite[Proposition~3]{Sanz2010Quantum} identifies strong irreducibility
    with eventual full Kraus rank, hence with normality.
\end{proof}
